\chapter{Related Work}

This chapter reviews prior work related to time-series anomaly detection and concept drift adaptation. Existing subsequence anomaly detection methods provide strong foundations for detecting unusual time-series regions. In particular, this chapter focuses on subsequence anomaly detection methods used as baselines in the AnDri work, including NormA, SAND, and DAMP, and then discusses concept drift detection methods and benchmark resources. 

\section{Normal-Model-Based Detection: NormA}

NormA is an unsupervised subsequence anomaly detection algorithm for time-series data \citep{boniol2021norma}. Instead of detecting only isolated discords, NormA builds a normal model from frequently occurring subsequences and measures anomalies based on their distance from this model. The method is scalable and does not require labeled data.

However, NormA mainly relies on a global normal model learned from historical data. In environments with concept drift, normal behavior may change over time, making it difficult to distinguish between true anomalies and newly emerging normal patterns. AnDri addresses this limitation by dynamically managing active and inactive normal patterns \citep{park2025adaptive}.

\section{Streaming Detection: SAND and DAMP}

SAND extends subsequence anomaly detection to streaming data \citep{boniol2021sand}. It incrementally updates its model as new data arrives, which makes it suitable for real-time monitoring scenarios such as sensors and system logs. However, SAND mainly focuses on automatic anomaly detection and does not explicitly support user exploration of concept drift and changing normal behavior.

DAMP is also a streaming anomaly detection method based on Matrix Profile \citep{lu2022damp}. Its main advantage is computational efficiency, allowing anomaly detection on very large and fast data streams. DAMP is effective for discord discovery, but it does not explicitly model concept drift. As normal patterns evolve or recur over time, discord-based methods may struggle to separate long-term drift from short-term anomalies \citep{park2025adaptive}.

Compared with these methods, AnDri focuses not only on anomaly detection performance but also on dynamically modeling changing normal behavior \citep{park2025andri}.



\section{Concept Drift Detection}

Concept drift refers to data distribution changes over time. Drift adaptation is a central problem in evolving data streams \citep{gama2014survey,lu2019learning}. ADWIN uses an adaptive sliding window to detect distribution changes and adjust the window size according to observed evidence of drift \citep{bifet2007adwin}. In real time series, anomaly and concept drift can look similar: both appear as deviations from past behavior. If a system treats all change as anomalous, it may produce many false positives during drift. If it adapts too aggressively, it may absorb true anomalies into the normal model \citep{park2025andri,park2025adaptive}.